\documentclass[12pt,a4paper]{article}
\usepackage[utf8]{inputenc}
\usepackage[T1]{fontenc}
\usepackage[spanish,es-tabla]{babel}
\usepackage{lmodern}
\usepackage[margin=2.5cm]{geometry}
\usepackage{graphicx}
\usepackage{float}
\usepackage{booktabs}
\usepackage{tabularx}
\usepackage{array}
\usepackage{enumitem}
\usepackage{xcolor}
\usepackage{tikz}
\usetikzlibrary{positioning,arrows.meta,calc}
\usepackage{hyperref}
\usepackage{setspace}
\usepackage{fancyhdr}
\setlength{\headheight}{15pt}
\usepackage{pdflscape}
\graphicspath{{../diagramas/}}
\hypersetup{colorlinks=true,linkcolor=blue!50!black,urlcolor=blue!50!black}
\makeatletter\g@addto@macro\UrlBreaks{\do\-}\makeatother
\setstretch{1.1}
\setlength{\parindent}{0pt}
\setlength{\parskip}{6pt}
\newcommand{\pendiente}[1]{\textcolor{red}{\textbf{[POR COMPLETAR: #1]}}}
\pagestyle{fancy}\fancyhf{}
\lhead{Banco de Preguntas Saber}
\rhead{Documento de arquitectura -- Segunda iteración}
\cfoot{\thepage}
\begin{document}
\begin{titlepage}
\centering
{\Large Universidad del Cauca\\Facultad de Ingeniería Electrónica y Telecomunicaciones\\Programa de Ingeniería de Sistemas\par}
\vspace{2cm}
{\huge\bfseries Banco de Preguntas Saber\par}
\vspace{0.5cm}
{\Large Documento de Arquitectura -- Iteración 2\par}
\vspace{1cm}
{Equipo ``Sin Nombre''\par}
\vspace{0.4cm}
Juan Manuel Meneses Daza\\Juan Sebastian Muñoz Ruiz\\Juan Sebastian Molano Bolaños\\Santiago Rodríguez Bolaños
\vfill
Ingeniería de Software II -- 2026.2\\Popayán, noviembre de 2026
\end{titlepage}
\tableofcontents\newpage

\section{Resumen}
Este documento describe lo verificado en el repositorio del Banco de Preguntas Saber Pro. La iteración separa autenticación, banco de preguntas, revisión por pares y auditoría en servicios independientes. La interfaz de escritorio conserva operaciones locales del primer corte; por ello la migración no es completa. La escalabilidad operativa no se ha medido ni demostrado. La base principal de esta descripción es el código de los cuatro servicios, \texttt{docs/decisiones-iteracion-2.md} y \texttt{entregables/sprint-2/SAD\_Banco\_Preguntas\_Iteracion2.pdf}.

\section{Requisitos funcionales}
Las historias y criterios se expresan en lenguaje de negocio. El alcance de HU01--HU04 procede de \texttt{docs/decisiones-iteracion-2.md}; la parte de HU05 se limita a la revisión implementada.

\subsection{Épica: administrar usuarios (HU01)}
\textbf{Historia.} Como persona usuaria, quiero registrarme e iniciar sesión para acceder a las funciones que corresponden a mi rol.

\textbf{Criterios de aceptación observables}
\begin{enumerate}[leftmargin=*]
\item Al registrar un login nuevo con datos válidos, el sistema confirma el alta y la cuenta queda activa.
\item Al intentar registrar un login ya existente, el sistema rechaza el alta sin crear una segunda cuenta.
\item El registro rechaza una contraseña que no alcance seis caracteres o no incluya un dígito, una mayúscula y un carácter especial.
\item La contraseña no aparece en la respuesta de alta ni en el listado de usuarios; se almacena protegida.
\item Al iniciar sesión con credenciales válidas de una cuenta activa, el sistema devuelve una credencial temporal firmada con la que los servicios autorizados aceptan las operaciones permitidas.
\item Una cuenta inactiva o credenciales incorrectas no permiten iniciar sesión y reciben una respuesta de acceso denegado.
\end{enumerate}

\subsection{Épica: gestionar preguntas (HU02)}
\textbf{Historia.} Como autor, quiero crear, consultar, filtrar y corregir preguntas para construir un banco útil para la preparación Saber Pro.

\textbf{Criterios de aceptación observables}
\begin{enumerate}[leftmargin=*]
\item Un autor o administrador puede registrar una pregunta con sus datos académicos, contexto, pregunta directa y opciones; si es válida, queda guardada en estado Borrador.
\item Al consultar una pregunta existente, el sistema muestra su contenido, respuesta correcta y estado; una pregunta inexistente se informa como no encontrada.
\item La búsqueda puede limitar resultados por estado, dificultad, competencia, tema, subtema o texto, de acuerdo con los filtros solicitados.
\item El autor o administrador puede corregir el contenido mientras la pregunta está en Borrador.
\item Si se intenta modificar el contenido fuera de Borrador, el sistema rechaza la operación y conserva el contenido previamente guardado.
\item Un usuario sin token válido no puede acceder a las operaciones del Banco; una operación de escritura solo la realizan los roles autorizados.
\end{enumerate}

\subsection{Épica: asegurar calidad de preguntas (HU03)}
\textbf{Historia.} Como autor, quiero recibir una explicación cuando una pregunta no cumple las reglas acordadas para corregirla antes de que se incorpore al banco.

\textbf{Alcance funcional.} El servicio valida la estructura al crear o modificar una pregunta y devuelve el resultado de validación; no existe un paso de revisión humana para estas reglas.

\textbf{Criterios de aceptación observables}
\begin{enumerate}[leftmargin=*]
\item El sistema rechaza un contexto ausente o menor de veinte caracteres.
\item El sistema rechaza una pregunta directa que no cumpla el formato acordado: longitud mínima, un único signo de interrogación y signo final de pregunta.
\item El sistema exige cinco opciones de respuesta en total, una correcta y cuatro distractores.
\item El sistema rechaza las expresiones «Todas las anteriores» y «Ninguna de las anteriores» en las opciones.
\item Cada opción debe tener contenido, al menos tres caracteres, iniciar con mayúscula y no duplicar otra opción.
\item El índice de respuesta correcta debe señalar una opción existente y no vacía, sin duplicidad con otra opción.
\item Nombre, justificación, bibliografía, competencia, tema, subtema y dificultad deben estar presentes.
\item Cuando una regla falla, el sistema identifica el problema y no guarda la pregunta como válida; cuando todas se cumplen, permite continuar con la operación solicitada.
\end{enumerate}

\subsection{Épica: administrar el ciclo de vida (HU04)}
\textbf{Historia.} Como autor, revisor o administrador, quiero que cada pregunta avance solo por etapas permitidas y que el historial muestre quién realizó cada cambio.

\textbf{Flujo de estados.} Borrador → Pendiente de revisión → En revisión → Aprobada → Publicada → Archivada. También se admiten los retornos y rechazos definidos por el código: Pendiente de revisión puede volver a Borrador; En revisión puede terminar en Aprobada o Rechazada; Rechazada puede volver a Borrador o archivarse; Aprobada puede archivarse. Archivada es terminal.

\textbf{Criterios de aceptación observables}
\begin{enumerate}[leftmargin=*]
\item Una pregunta nueva comienza en Borrador y el sistema permite únicamente las transiciones definidas para su estado actual.
\item El sistema rechaza una transición no permitida y conserva el estado vigente.
\item Todo cambio de estado efectivo queda registrado con estado anterior, estado nuevo, usuario y fecha, y puede consultarse en Auditoría.
\item Solicitar el estado que la pregunta ya tiene no crea un segundo registro de cambio.
\item Una pregunta archivada no puede reactivarse ni borrarse físicamente.
\item Crear o editar contenido corresponde a Administrador o Autor; cambiar estado corresponde a Administrador, Autor o Revisor, siempre con credencial válida.
\end{enumerate}

\subsection{HU05 implementada parcialmente: revisión por pares}
\textbf{Historia.} Como revisor, quiero tomar una revisión y registrar un dictamen para que el Banco avance la pregunta.

\textbf{Criterios de aceptación observables}
\begin{enumerate}[leftmargin=*]
\item Cuando una pregunta llega a Pendiente de revisión, el servicio de Revisión crea una revisión pendiente a partir del evento recibido.
\item El autor que solicitó la revisión no puede revisar su propia pregunta.
\item Solo un revisor autorizado puede tomar una revisión pendiente; una revisión tomada no puede ser tomada por otro revisor.
\item Al tomarla, el servicio solicita al Banco que avance la pregunta a En revisión; si el Banco no confirma, la revisión no se presenta como tomada correctamente.
\item El revisor que tomó la revisión puede emitir un dictamen de aprobación o rechazo; el Banco realiza el cambio de estado.
\item Un rechazo sin observaciones se rechaza.
\item Las entregas repetidas del mismo evento no crean revisiones duplicadas.
\end{enumerate}

\textbf{Límite de alcance.} La asignación automática de revisores y una rúbrica/formato completo de evaluación no están implementados; por ello HU05 se considera parcial, no aceptada en su totalidad.

\subsection{Fuente de verificación}
Los criterios anteriores resumen el comportamiento comprobado en los controladores, servicios, filtros y pruebas automatizadas de los cuatro módulos; las reglas detalladas se mantienen en \texttt{RULES.md}. Constituyen criterios funcionales comprobables para esta versión. El registro del Product Backlog y la conformidad formal del Product Owner son evidencias de gestión y no se infieren del código.

\subsection{Roles y permisos observados}
\begin{tabularx}{\textwidth}{p{3.5cm}X}
\toprule Operación & Permiso verificado \\
\midrule
Crear o editar pregunta & \texttt{ADMINISTRADOR} o \texttt{AUTOR\_PREGUNTAS} \\
Cambiar estado & \texttt{ADMINISTRADOR}, \texttt{AUTOR\_PREGUNTAS} o \texttt{REVISOR} \\
Consultar preguntas y estadísticas & Se valida token; el controlador no aplica un rol específico \\
Tomar revisión o emitir dictamen & Revisor o Administrador, con reglas de pertenencia \\
\bottomrule
\end{tabularx}

\section{Prototipos y usabilidad}
En los artefactos revisados no se encontró evidencia de prototipos ni de una evaluación de usabilidad correspondiente a esta arquitectura. La usabilidad se trata mediante el escenario de calidad propuesto en la sección siguiente.

\section{Escenarios de calidad}
Los escenarios técnicos del SAD y MiniQAW se reexpresan como esfuerzo de las personas involucradas, no como rendimiento medido. Todas las cifras siguientes son \textbf{propuesta, por validar con el equipo y el profesor}; no son resultados obtenidos.
\begin{tabularx}{\textwidth}{p{3cm}X p{3cm}}
\toprule Atributo & Escenario de negocio & Medida propuesta \\
\midrule
Modificabilidad & Agregar o ajustar una regla de calidad de pregunta y revisar que no se afecten las demás reglas. & 4 personas-hora para cambio y pruebas; propuesta, por validar. \\
Seguridad & Revisar roles, token y respuestas ante intentos de acceso indebido. & 6 personas-hora para revisión manual y casos de prueba; propuesta, por validar. \\
Usabilidad & Una persona nueva crea y busca preguntas con instrucciones mínimas. & 30 minutos de persona usuaria para completar la primera tarea; propuesta, por validar. \\
Escalabilidad & Preparar una demostración de mayor carga y registrar el resultado. & 8 personas-hora para preparar y ejecutar una prueba inicial; propuesta, por validar. No se fija umbral de usuarios porque el código no tiene medición ni infraestructura escalable verificada. \\
Trazabilidad & Revisar si una pregunta que cambia de estado queda registrada en auditoría y revisión. & 4 personas-hora para prueba y análisis de fallos; propuesta, por validar. \\
\bottomrule
\end{tabularx}

\section{Modelo de datos}
Las tablas se verifican en las clases de configuración de Auth, Banco, Revisión y Auditoría. Los catálogos de roles y estados se almacenan como texto, no como tablas independientes. Las referencias de pregunta desde revisión/auditoría no son llaves foráneas entre bases.
\begin{figure}[p]\centering\includegraphics[width=0.78\textwidth]{06-er-bases.png}\caption{Esquemas por base de datos.}\end{figure}

\section{Arquitectura}
\subsection{Vista de contexto}
\begin{figure}[p]\centering\includegraphics[width=0.78\textwidth]{01-contexto-c4.png}\caption{C4 nivel 1.}\end{figure}
\subsection{Vista de contenedores}
\begin{landscape}\begin{figure}[H]\centering\includegraphics[width=0.94\linewidth]{02-contenedores-c4.png}\caption{C4 nivel 2: cuatro servicios, broker, bases y cliente.}\end{figure}\end{landscape}
\subsection{Vista de componentes}
\clearpage
\begin{landscape}\begin{figure}[p]\centering\includegraphics[width=0.82\linewidth]{03-componentes-c4.png}\caption{C4 nivel 3: Banco, Revisión y Auditoría.}\end{figure}\end{landscape}
\subsection{Vistas 4+1}
\begin{itemize}[leftmargin=*]
\item Lógica: dominios de preguntas, usuario, revisión y auditoría; véanse diagramas UML.
\item Desarrollo: cuatro módulos Maven y cliente JavaFX heredado.
\item Procesos: REST para operaciones y RabbitMQ para eventos; los consumidores registran identificadores procesados.
\item Física: configuración local con SQLite por servicio; `docker-compose.yml` define solo el broker.
\item Escenarios: crear pregunta y revisión por pares, descritos más adelante.
\end{itemize}
\subsection{Comunicación y eventos}
Auth emite tokens; Banco, Revisión y Auditoría los verifican localmente. Banco publica `question.created` y `question.state-changed` en `saberpro.events`; Auditoría consume ambos por `auditoria.preguntas`, Revisión consume cambios de estado por `revision.preguntas`. Revisión llama al Banco con el token del revisor. La publicación no tiene Outbox y los reintentos de consumidores no sustituyen a una DLQ.
\begin{figure}[p]\centering\includegraphics[width=0.78\textwidth]{07-secuencia-crear-pregunta.png}\caption{Secuencia de creación.}\end{figure}
\begin{figure}[p]\centering\includegraphics[width=0.78\textwidth]{08-secuencia-revision-pares.png}\caption{Secuencia de revisión por pares.}\end{figure}
\begin{figure}[p]\centering\includegraphics[width=0.78\textwidth]{09-estados.png}\caption{Estados de pregunta.}\end{figure}

\section{Información repetida entre servicios}
Se repiten roles, formato de token, códigos de estado y contratos de eventos. Se conservan copias locales para que los servicios no dependan de clases de otro módulo. Los cambios requieren compatibilidad y pruebas coordinadas; las referencias \texttt{question\_id} entre bases no son llaves foráneas. Ver \texttt{docs/decisiones-iteracion-2.md}.

\section{Decisiones finales y patrones aplicados}
Decisiones finales: microservicios con eventos; una SQLite por servicio; token HMAC-SHA256 compartido; validación embebida; \texttt{revision-service} en lugar del servicio anterior de validación; jerarquía State; Pipeline de filtros; consumidores idempotentes por \texttt{eventId}. Los ADR individuales están en \texttt{docs/adr/}. Observer permanece en el cliente local de \texttt{src/}; no es el mecanismo de eventos entre servicios. La lista académica de patrones «vistos» requiere validación del profesor.

\section{Cumplimiento del Sprint 2}
La matriz detallada de evidencia y acciones está en \texttt{docs/CUMPLIMIENTO\_SPRINT2.md}. Conclusión: cumplimiento parcial; escalabilidad no demostrada.

\section{Pruebas}
\begin{itemize}
\item \texttt{auth-service}: 21; pasó.
\item \texttt{banco-service}: 94; pasó.
\item \texttt{revision-service}: 40; pasó.
\item \texttt{auditoria-service}: 9; pasó.
\item Monolito raíz: 158; pasó.
\end{itemize}
Ejecutadas con Maven el 2026-11-22. Los conteos no equivalen a cobertura; no se encontró medición de cobertura.

\section{Gestión del proyecto}
El Product Backlog, Sprint Backlog, Task Board y Burndown Chart se gestionan en Jira/Trello y se entregan mediante el acceso al proyecto. No se duplican como artefactos generados en este SAD.

\section{Sustentación y repositorio}
Repositorio de código: \href{https://github.com/BastianRu/proyecto-saberpro-ing-software-II}{GitHub del proyecto}. La participación de los integrantes se consulta en el historial del repositorio. La evidencia audiovisual se entrega por el medio definido para el curso.

\end{document}
